\documentclass[12pt, a4paper]{article}

% Paquetes esenciales para español, matemáticas y gráficos
\usepackage[utf8]{inputenc}
\usepackage[spanish, es-tabla]{babel}
\usepackage{amsmath}
\usepackage{graphicx}
\usepackage{hyperref}
\usepackage{geometry}
\geometry{a4paper, margin=2.5cm}

% Información del documento
\title{\textbf{Análisis Técnico y Evaluación del Modelo de Síntesis FM}}
\author{Tu  Nombre \\ \textit{Borrador para reunión con experto}}
\date{\today}

\begin{document}

\maketitle

\begin{abstract}
Este documento detalla el proceso de evaluación del modelo de predicción de parámetros para síntesis FM. Se analiza la discrepancia entre métricas de clasificación global (FAD) y la percepción auditiva del tono (Pitch), proponiendo un sistema de validación basado en búsqueda en rejilla (Grid Search) utilizando coeficientes cepstrales (MFCC) para solventar los problemas de no-unicidad armónica.

\end{abstract}

\tableofcontents
\newpage

\section{Evaluación y Métricas del Modelo}
% Aquí introduciremos por qué es difícil evaluar la síntesis FM.


\subsection{Limitaciones de Fréchet Audio Distance (FAD)}
% Aquí redactaremos cómo sacamos un 0.99 de FAD y por qué engaña con el Pitch.
Para la evaluación del modelo, se seleccionó la métrica \textit{Fréchet Audio Distance} (FAD) \cite{kilgour2019}. 
Esta herramienta, introducida originalmente por investigadores de Google AI, es utilizada actualmente en proyectos muy grandes como 
MusicLM \cite{agostinelli2023}, pudiendo considerarse como el \textbf{estado del arte} para medir la calidad objetiva de los audios generados mediante Inteligencia Artificial, adaptando al dominio del audio las métricas utilizadas en la generación de imágenes.

Conceptualmente, el FAD deriva de \textit{Fréchet Inception Distance} (FID), una métrica referente en la evaluación de modelos generativos de imágenes.
Al igual que su predecesora, esta métrica no realiza comparación muestra a muestra, sino que compara características de alto nivel extraídas por una red neuronal profunda.

Para la extracción de estas caractarísticas, el FAD utiliza un modelo llamado \textit{VGGish}, un clasificador de audio entrenado con una base de datos 
de gran envergadura de vídeos de YouTube. Este modelo transforma la señal de audio en \textit{log-mel features} extraídas del espectrograma de magnitud

Finalmente, la métrica se obtiene calculando la \textit{Distancia de Fréchet} entre la música de referencia y la generada. Un valor de FAD reducido indíca una estadística similar entre el audio de referencia y el generado.


Para la implementación de la métrica FAD en el proyecto, se compararon dos conjuntos de datos: una muestra de 2.000 audios FM generados aleatoriamente
y sus predicciones con el modelo utilizado en el \textit{Prototipo 4}. El resultado obtenido fue de \textbf{0.9999}, un valor que, según la literatura de referencia, indica una similitud estadística casi absoluta entre las distribuciones de audio.


A pesar de que el resultado obtenido planteaba una evaluación teóricamente óptima, el análisis cualitativo reveló una discrepancia. 
El motivo de la poca validez de la métrica en el proyecto, se atribuye a un sesgo de interpretación de su 
utilidad real, ya que el FAD está diseñado para evaluar la calidad global de modelos generativos \cite{kilgour2019}, mientras que este proyecto se centra 
en la estimación de los parámetros que luego utiliza para generar audio FM. 
Puesto que tanto el conjunto de referencia como las predicciones se crean mediante síntesis FM, su huella espectral y sus caractarísticas de alto nivel(\textit{embeddings})
son casi idénticas \cite{kilgour2019}. Por tanto, el FAD evalúa positivamente que el modelo ha aprendido a recrear el timbre o "textura" de 
la síntesis FM, pero, como se comentaba previamente, resulta una muestra ciega a precisión muestra a muestra. En resumen, se prioriza la 
verosimilitud estadística del sonido sobre la identidad exacta del mismo \cite{kilgour2019, agostinelli2023}.


\subsection{Validación por Fuerza Bruta (Grid Search)}
% Explicaremos la necesidad de buscar en todo el espacio de parámetros.
Una de las cuestiones principales que se plantean en el desarrollo de este proyecto es si el uso de un modelo de Inteligencia Artificial representa
la aproximación más óptima para el objetivo de la inferencia de parámetros. Esta disyuntiva surge, sobre todo, al analizar los primeros prototipos, 
como el \textit{Prototipo 4}, en el que se emplean únicamente 3 parámetros para la síntesis FM: \textit{Carrier}, \textit{Index} y \textit{Ratio}.
Debido al reducido tamaño de este espacio de parámetros, resulta lógico plantear en la viabilidad de realizar un método de búsqueda exhaustiva 
(\textit{Grid Search}). Este enfoque consiste en realizar un recorrido iterativo por todas las combinaciones posibles por ellos que prueben 
el total de combinaciones posibles, evaluándolas y registrando aquellas configuraciones que presenten una mayor similitud acústica con
respecto al sonido objetivo.

Por otro lado, surge la necesidad de encontrar un método rubusto para evaluar la fidelidad del modelo de inferencia de parámetros, buscando una métrica
que simule la percepción de un oído humano experto haría un oído humano experto. Esta necesidad y la cuestión anterior convergen en la fase de experimentación.
De este modo, se consiguen probar diferentes métricas que evalúen la similitud acústica dentro del algoritmo de \textit{Grid Search}, comprobando
simultáneamente la viabilidad técnica del enfoque de fuerza bruta frente al modelo predictivo.

Para lograr esta simulación del sistema auditivo humano, en el procesamiento digital de señales de audio se encuentra la \textbf{Escala Mel} como estándar.
Esta aplica una transformación logarítmica que modela la resol


\subsubsection{Comparativa: Espectrograma de Mel vs. MFCC}
% Aquí meteremos los resultados que sacamos hoy: cómo Mel falla por la rejilla armónica y MFCC acierta capturando la envolvente espectral.

\section{Evolución del Modelo: Prototipo 5}
% Esta es la parte de tu compañero que vas a estudiar e ir rellenando poco a poco.
\subsection{Ampliación del espacio de parámetros: La Envolvente}
% Explicación de los nuevos parámetros para hacer el sonido más "orgánico".

\vspace{2cm} % Esto añade un poco de espacio en blanco
\begin{thebibliography}{9}

\bibitem{kilgour2019}
Kilgour, K., Zuluaga, M., Roblek, D., \& Sharifi, M. (2019). 
\textit{Fréchet Audio Distance: A Reference-free Metric for Evaluating Music Enhancement Algorithms}. 
Interspeech 2019, Graz, Austria, pp. 2350-2354.

\bibitem{agostinelli2023}
Agostinelli, A., Denk, T. I., Borsos, Z., Engel, J., Verzetti, M., \& Roberts, A. (2023). 
\textit{MusicLM: Generating Music From Text}. 
arXiv preprint arXiv:2301.11325.

\end{thebibliography}
\end{document}